% TOP_PROBLEM: {"id":30007697,"problem_number":"LOCAL-30007697","title":"Training-data compensation","category_id":21,"category":{"id":21,"name":"economics","display_name":"Economics","description":"Open economic theory statements, published research agendas, and proposed scoped specifications; question provenance and historical priority are qualified per record.","slug":"economics","order_index":21,"created_at":"2026-10-08T00:00:00Z"},"status":"provenance_pending","external_url":null,"legacy_ids":[],"published":false,"statement":"Design feasible training-content compensation mechanisms balancing model value, original-content quality, creator participation and developer incentives.","economics":{"original_id":186,"field":"Innovation and AI economics","kind":"design","formalization_basis":"proposed_specification","source_status":"not_verified","priority_status":"not_established","priority_note":"No exact open-problem passage was verified. The supporting publication does not establish priority or unresolved status.","earliest_verified_reference":null,"current_reference":{"authors":["Gaétan de Rassenfosse","Adam B. Jaffe","Joel Waldfogel"],"title":"Intellectual Property and Creative Machines","year":2024,"url":"https://www.nber.org/papers/w32698","doi":"10.3386/w32698","locator":"Primary abstract and indexed chapter conclusion snippets; complete agenda paragraph not retrieved","evidence_summary":"Addresses economic IP trade-offs for creative machines; exact compensation and output-protection agenda passages remain unverified."},"formalization_note":"The cited primary work supports the subject or supplies solutions in particular settings, but no exact open-question passage for this specification was verified. The mathematical target and restrictions are proposed here, with no canonical-conjecture or priority claim."},"statusQualification":"Provenance pending: an explicit published statement of this exact open question was not verified. This is a catalogue-authored candidate specification, not a certified open conjecture. The cited primary work supports the subject or supplies solutions in particular settings, but no exact open-question passage for this specification was verified. The mathematical target and restrictions are proposed here, with no canonical-conjecture or priority claim.","statusReviewedAt":"2026-10-08"}
% FIRST_POSTED: 2026-10-08
% ENTRY_KIND: statement_only
% !TeX program = lualatex
\documentclass[11pt]{article}
\usepackage{fontspec}
\usepackage[a4paper,margin=25mm]{geometry}
\usepackage{amsmath,amssymb,mathtools}
\usepackage{xurl}
\usepackage[hidelinks]{hyperref}
\newcommand{\sref}[2]{\hyperlink{source:#1:#2}{[S#2]}}
\setlength{\emergencystretch}{3em}
\begin{document}
% problemId:30007697
\section*{Training-data compensation}\label{30007697}
\subsection{Definitions and mathematical statement}
Fix a market containing original-content creators, one model developer and downstream users. Creator \(i\) chooses costly original-content quality \(x_i\geq0\), with real cost \(c_i(x_i)\); the developer obtains licensed training content and produces model quality \(q(x)\), where \(x\) is the vector of creator choices. A compensation mechanism \(r\) specifies admissible training uses, observable quality verification and transfers \(t_i(x)\) to creators. Let \(\mathcal R\) be a stated feasible mechanism class satisfying budget balance, creator participation and the developer's participation constraint. Creator utility is \(t_i(x)-c_i(x_i)\), plus separately specified revenue from nontraining uses. Let \(B(q,x)\) be users' aggregate gross value from the model and original works, and \(K(q)\) the developer's real training and deployment cost. With transfers counted only once, social welfare is
\[
W(r)=B(q(x(r)),x(r))-\sum_i c_i(x_i(r))-K(q(x(r))),
\]
where \(x(r)\) is an equilibrium quality profile under mechanism \(r\). Determine when a welfare-maximizing \(r^\ast\in\arg\max_{r\in\mathcal R}W(r)\) can be implemented using verifiable usage and quality information, and how it changes future content supply. This is a proposed mechanism problem; conclusions depend on specified outside options, production technology and equilibrium selection, and assert no universal legal compensation rule.

\par\textbf{Formulation and scope.} The cited primary work supports the subject or supplies solutions in particular settings, but no exact open-question passage for this specification was verified. The mathematical target and restrictions are proposed here, with no canonical-conjecture or priority claim.

\par\textbf{Open-question provenance.} An explicit published open-question statement was not verified. Sources below are supporting literature and must not be treated as a first listing.

\par\textbf{Historical priority.} No exact open-problem passage was verified. The supporting publication does not establish priority or unresolved status.

\subsection{Short English statement}
Design feasible training-content compensation mechanisms balancing model value, original-content quality, creator participation and developer incentives.

\subsection{Sources}
\begin{itemize}\raggedright
\item[\textnormal{[S1]}]\hypertarget{source:30007697:1}{} Gaétan de Rassenfosse, Adam B. Jaffe, Joel Waldfogel. 2024. Intellectual Property and Creative Machines. Primary abstract and indexed chapter conclusion snippets; complete agenda paragraph not retrieved.
\url{https://www.nber.org/papers/w32698}.
DOI: 10.3386/w32698.
{\small\textit{Supporting or current literature; see evidence qualification. Addresses economic IP trade-offs for creative machines; exact compensation and output-protection agenda passages remain unverified.}}
\end{itemize}
\end{document}
